\documentclass[12pt,a4paper]{article}
\usepackage[utf8]{inputenc}
\usepackage[T1]{fontenc}
\usepackage[french]{babel}
\usepackage{amsmath,amssymb,amsthm}
\usepackage{graphicx}
\usepackage{tikz}
\usepackage{pgfplots}
\pgfplotsset{compat=1.18}
\usepackage{booktabs}
\usepackage{array}
\usepackage{multirow}
\usepackage{geometry}
\usepackage{hyperref}
\usepackage{listings}
\usepackage{xcolor}
\usepackage{float}
\usepackage{caption}
\usepackage{subcaption}
\usepackage{fancyhdr}
\usepackage{enumitem}

% Configuration de la page
\geometry{margin=2.5cm}
\pagestyle{fancy}
\fancyhf{}
\fancyhead[L]{\leftmark}
\fancyhead[R]{HIGHLIGHT+}
\fancyfoot[C]{\thepage}

% Configuration des hyperliens
\hypersetup{
    colorlinks=true,
    linkcolor=blue,
    filecolor=magenta,      
    urlcolor=cyan,
    pdftitle={Rapport HIGHLIGHT+},
    pdfauthor={Équipe HIGHLIGHT+}
}

% Commandes personnalisées
\newcommand{\highlight}{\textbf{HIGHLIGHT+}}
\newcommand{\teacher}{\textsc{Teacher}}
\newcommand{\student}{\textsc{Student}}
\newcommand{\gp}{\textsc{GP}}
\newcommand{\rl}{\textsc{RL}}

% Configuration des listes
\setlist[itemize]{leftmargin=*,topsep=0pt}
\setlist[enumerate]{leftmargin=*,topsep=0pt}

\title{\textbf{Rapport de Présentation Détaillé}\\[0.5em]
\Large Système de Détection Intelligente de Micro-Fuites de Méthane\\[0.3em]
\large Concours Innovation Natran x Fondation UTT - 2025}
\author{Équipe HIGHLIGHT+\\Housséni YABRE, Kabinet SYLLA, Nobert Bassooma DIDANERA}
\date{\today}

\begin{document}

\maketitle
\thispagestyle{empty}

\newpage
\tableofcontents
\newpage

\section{Introduction et Contexte}

\subsection{Problématique}

La détection de micro-fuites de méthane représente un défi majeur pour l'industrie énergétique. Les méthodes traditionnelles (patrouilles systématiques, capteurs fixes) présentent plusieurs limitations :

\begin{itemize}
    \item \textbf{Coûteuses} en énergie et temps
    \item \textbf{Peu efficaces} avec un taux de détection inférieur à 15\%
    \item \textbf{Non adaptatives} aux conditions environnementales variables
\end{itemize}

\subsection{Solution Proposée : \highlight}

\highlight{} est un système d'intelligence artificielle qui transforme un drone-dirigeable en détective autonome de micro-fuites de méthane. L'innovation réside dans l'\textbf{architecture Teacher-Student} combinant :

\begin{itemize}
    \item \textbf{Apprentissage actif} (Processus Gaussiens) pour la planification stratégique
    \item \textbf{Apprentissage par renforcement} pour la navigation tactique optimale
    \item \textbf{Distillation de connaissance} pour transférer l'expertise du \teacher{} au \student{}
\end{itemize}

\subsection{Objectifs}

Les objectifs principaux du système sont :

\begin{itemize}
    \item \textbf{Maximiser le taux de détection} : Objectif supérieur à 90\%
    \item \textbf{Minimiser la consommation énergétique} : Optimisation des trajectoires
    \item \textbf{Localisation précise} : Erreur inférieure à 2 mètres
    \item \textbf{Autonomie complète} : Décisions en temps réel sans intervention humaine
\end{itemize}

\subsection{Explication des Concepts Clés}

Pour mieux comprendre le fonctionnement de \highlight{}, voici une explication simplifiée des concepts fondamentaux utilisés dans le système.

\paragraph{Qu'est-ce qu'un Processus Gaussien (\gp{}) ?}

Imaginez que vous cherchez une source de méthane dans une grande zone. À chaque endroit où le drone mesure la concentration, vous obtenez une valeur. Un \textbf{Processus Gaussien} est un outil mathématique qui permet de :
\begin{itemize}
    \item \textbf{Prédire} la concentration de méthane en n'importe quel point de la zone, même là où le drone n'est pas encore allé
    \item \textbf{Estimer l'incertitude} de cette prédiction (à quel point on est sûr ou pas sûr de la valeur)
    \item \textbf{Créer une carte} de concentration qui s'améliore au fur et à mesure que le drone collecte plus de mesures
\end{itemize}

En d'autres termes, le \gp{} apprend des mesures déjà prises pour deviner ce qui se passe ailleurs, tout en indiquant à quel point ces prédictions sont fiables. C'est comme un détective qui utilise les indices collectés pour reconstituer une image complète de la scène.

\paragraph{Architecture Teacher-Student : Comment ça fonctionne ?}

Le système \highlight{} utilise une approche inspirée de l'éducation, avec deux composants qui travaillent ensemble :

\begin{itemize}
    \item \textbf{Le \teacher{} (\gp{})} : C'est l'\textit{expert} qui connaît bien la physique des panaches de méthane. Il analyse toutes les mesures collectées et décide \textit{stratégiquement} où le drone devrait aller ensuite pour maximiser les chances de trouver la source. Il pense à long terme et planifie l'exploration de manière intelligente.
    
    \item \textbf{Le \student{} (\rl{})} : C'est l'\textit{apprenti} qui apprend à naviguer efficacement. Au début, il ne sait pas grand-chose, mais en observant le \teacher{} et en pratiquant, il apprend progressivement les meilleures trajectoires. Il devient de plus en plus rapide et efficace, comme un pilote qui gagne en expérience.
    
    \item \textbf{La Distillation de Connaissance} : C'est le processus par lequel le \student{} apprend du \teacher{}. Au lieu de réinventer la roue, le \student{} copie les bonnes stratégies du \teacher{}, ce qui accélère son apprentissage. C'est comme un étudiant qui apprend d'un professeur expérimenté plutôt que de tout découvrir par lui-même.
\end{itemize}

\paragraph{Pourquoi cette combinaison est-elle puissante ?}

Le \teacher{} est excellent pour la \textbf{planification stratégique} : il sait où chercher et pourquoi. Cependant, il peut être un peu lent car il doit calculer beaucoup de choses à chaque étape. Le \student{}, une fois entraîné, est excellent pour la \textbf{navigation tactique} : il peut réagir rapidement et suivre des trajectoires optimisées. En les combinant, on obtient le meilleur des deux mondes : une exploration intelligente et une navigation rapide.

\paragraph{Comment le système détecte-t-il la fuite ?}

Le processus de détection se déroule en plusieurs étapes :

\begin{enumerate}
    \item \textbf{Collecte de mesures} : Le drone vole et mesure la concentration de méthane à différents endroits. Chaque mesure est comme un point de données sur une carte.
    
    \item \textbf{Modélisation avec \gp{}} : Le \teacher{} utilise toutes ces mesures pour construire une \textit{carte de probabilité} qui montre où la source est la plus susceptible de se trouver. Cette carte s'améliore à chaque nouvelle mesure.
    
    \item \textbf{Estimation de position} : Le système combine deux informations : (1) les zones où la concentration est la plus élevée, et (2) les zones où on est le plus confiant dans nos prédictions. La position estimée est là où ces deux critères sont optimaux.
    
    \item \textbf{Arrêt automatique} : Dès que le système est suffisamment confiant (plus de 85\% de certitude), il s'arrête automatiquement. Cela évite de gaspiller de l'énergie à continuer à chercher quand on a déjà trouvé la source.
\end{enumerate}

\paragraph{Qu'est-ce que l'Apprentissage par Renforcement (\rl{}) ?}

L'apprentissage par renforcement fonctionne sur le principe de la \textit{récompense} et de la \textit{pénalité}. Imaginez que vous apprenez à conduire :
\begin{itemize}
    \item Si vous prenez une bonne décision (par exemple, suivre une route qui mène rapidement à la source), vous recevez une \textbf{récompense positive}
    \item Si vous prenez une mauvaise décision (par exemple, voler dans une zone sans méthane), vous recevez une \textbf{pénalité} (consommation d'énergie inutile)
    \item Le \student{} apprend progressivement à maximiser les récompenses et minimiser les pénalités
    \item Au fil du temps, il développe une \textit{stratégie optimale} qui lui permet de trouver la source le plus rapidement possible tout en économisant l'énergie
\end{itemize}

\paragraph{En résumé : Comment tout fonctionne ensemble ?}

Le système \highlight{} fonctionne comme suit :
\begin{enumerate}
    \item Le drone commence à voler dans la zone de recherche
    \item Le \teacher{} analyse les mesures et suggère où aller ensuite (zones prometteuses ou zones incertaines à explorer)
    \item Le \student{} combine cette suggestion avec son expérience pour choisir la meilleure action immédiate
    \item Le drone se déplace et collecte une nouvelle mesure
    \item Le processus se répète, et à chaque itération, le système en sait plus sur la localisation de la source
    \item Quand la confiance atteint un seuil élevé, le système s'arrête et annonce la position estimée de la fuite
\end{enumerate}

Cette approche permet d'obtenir des résultats \textbf{excellents} (taux de détection supérieur à 90\%) tout en étant \textbf{efficace} (économie d'énergie de 35\%) et \textbf{rapide} (détection en moins de 2,5 secondes).

\section{Fondements Mathématiques}

\subsection{Modélisation Physique du Panache de Méthane}

\subsubsection{Équation d'Advection-Diffusion}

Le panache de méthane est modélisé par l'équation d'advection-diffusion :

\begin{equation}
\frac{\partial C}{\partial t} + \mathbf{u} \cdot \nabla C = D \nabla^2 C + Q \cdot \delta(x - x_0) \delta(y - y_0)
\end{equation}

où :
\begin{itemize}
    \item $C(x,y,t)$ : Concentration de méthane (kg/m³)
    \item $\mathbf{u} = (u_x, u_y)$ : Vecteur vitesse du vent (m/s)
    \item $D$ : Coefficient de diffusion (m²/s)
    \item $Q$ : Débit de la fuite (kg/s)
    \item $(x_0, y_0)$ : Position de la source
\end{itemize}

\subsubsection{Solution Analytique (Modèle Gaussien)}

Pour un panache en régime stationnaire avec diffusion gaussienne, la solution analytique est :

\begin{equation}
C(x,y,t) = \frac{Q}{2\pi \sigma_x \sigma_y u} \exp\left(-\frac{(x-x_0')^2}{2\sigma_x^2} - \frac{(y-y_0')^2}{2\sigma_y^2}\right) \exp(-\lambda t)
\end{equation}

où :
\begin{itemize}
    \item $\sigma_x, \sigma_y$ : Écarts-types de diffusion horizontale et verticale (m)
    \item $u = \|\mathbf{u}\|$ : Norme de la vitesse du vent (m/s)
    \item $(x_0', y_0') = (x_0 + u_x \cdot t, y_0 + u_y \cdot t)$ : Position effective (advection)
    \item $\lambda$ : Taux de décroissance temporelle (s⁻¹)
\end{itemize}

L'implémentation calcule cette concentration en trois étapes :
\begin{enumerate}
    \item Calcul de la position effective avec advection par le vent
    \item Calcul des distances relatives à la source
    \item Application du facteur de décroissance temporelle et de la distribution gaussienne
\end{enumerate}

\subsubsection{Gradient de Concentration}

Le gradient spatial est calculé analytiquement :

\begin{align}
\nabla C &= \left(\frac{\partial C}{\partial x}, \frac{\partial C}{\partial y}\right) \\
\frac{\partial C}{\partial x} &= -C(x,y) \times \frac{x - x_0'}{\sigma_x^2} \\
\frac{\partial C}{\partial y} &= -C(x,y) \times \frac{y - y_0'}{\sigma_y^2}
\end{align}

Le gradient indique la direction de la source et est utilisé pour guider la navigation du drone.

\subsection{Modèle de Processus de Décision Markovien (MDP)}

Le problème de détection est formulé comme un \textbf{MDP} $(S, A, P, R, \gamma)$ :

\subsubsection{Espace d'État $S$}

L'état $s \in S$ est un vecteur de 16 dimensions :

\begin{equation}
s = [x, y, z, v_x, v_y, v_z, \text{SNR}, \text{wind}_x, \text{wind}_y, \mu_{\text{GP}}, \sigma_{\text{GP}}, \text{grad}_x, \text{grad}_y, \text{energy}, \text{time}, \text{concentration}]
\end{equation}

où :
\begin{itemize}
    \item $(x, y, z)$ : Position du drone (m)
    \item $(v_x, v_y, v_z)$ : Vitesse du drone (m/s)
    \item $\text{SNR}$ : Rapport signal/bruit du capteur
    \item $(\text{wind}_x, \text{wind}_y)$ : Composantes du vent (m/s)
    \item $(\mu_{\text{GP}}, \sigma_{\text{GP}})$ : Prédiction et incertitude du \gp{}
    \item $(\text{grad}_x, \text{grad}_y)$ : Gradient de concentration
    \item $\text{energy}$ : Énergie consommée (J)
    \item $\text{time}$ : Temps écoulé (s)
    \item $\text{concentration}$ : Dernière concentration mesurée (kg/m³)
\end{itemize}

\subsubsection{Espace d'Action $A$}

L'action $a \in A$ est un vecteur de 3 dimensions :

\begin{equation}
a = [\Delta v, \Delta h, \Delta \psi]
\end{equation}

où :
\begin{itemize}
    \item $\Delta v$ : Variation de vitesse horizontale (m/s)
    \item $\Delta h$ : Variation d'altitude (m)
    \item $\Delta \psi$ : Variation de cap (rad)
\end{itemize}

Les actions sont normalisées dans $[-1, 1]$ et mappées aux commandes réelles du drone.

\subsubsection{Fonction de Transition $P$}

La transition d'état suit la dynamique du drone :

\begin{equation}
s_{t+1} = f(s_t, a_t) + \varepsilon
\end{equation}

où $f$ modélise :
\begin{itemize}
    \item La cinématique du drone (mouvement inertiel)
    \item L'effet du vent sur la trajectoire
    \item La consommation énergétique
\end{itemize}

\subsubsection{Fonction de Récompense Éco-Informative $R$}

La fonction de récompense combine deux objectifs :

\begin{equation}
R(s, a) = \alpha \cdot \Delta I(\mathcal{M}_{\text{GP}}) - \beta \cdot E(s, a)
\end{equation}

\paragraph{Gain d'Information $\Delta I(\mathcal{M}_{\text{GP}})$}

Le gain d'information mesure la réduction d'incertitude du modèle \gp{} :

\begin{equation}
\Delta I(\mathcal{M}_{\text{GP}}) = H(\mathcal{M}_{\text{GP}} | \text{observations}) - H(\mathcal{M}_{\text{GP}} | \text{observations} \cup \{x_{\text{new}}\})
\end{equation}

où $H$ est l'entropie (incertitude). En pratique :

\begin{equation}
\Delta I(\mathcal{M}_{\text{GP}}) \approx \frac{\sigma^2(x_{\text{new}})}{\sigma^2(x_{\text{new}}) + \sigma^2_{\text{noise}}}
\end{equation}

\paragraph{Coût Énergétique $E(s, a)$}

Le coût énergétique modélise la consommation du drone :

\begin{equation}
E(s, a) = c_1 \cdot \|a\|^2 + c_2 \cdot |\Delta h| + c_3 \cdot \|v_{\text{air}}\|^3
\end{equation}

où :
\begin{itemize}
    \item $\|a\|^2$ : Norme de l'action (effort de contrôle)
    \item $|\Delta h|$ : Variation d'altitude (coût vertical)
    \item $\|v_{\text{air}}\|^3$ : Vitesse relative à l'air au cube (résistance aérodynamique)
\end{itemize}

\paragraph{Poids d'Équilibrage $\alpha$ et $\beta$}

Les poids sont adaptatifs selon le contexte :

\begin{align}
\alpha &= \alpha_0 \times \text{SNR} \times (1 - \text{confidence}_{\text{GP}}) \\
\beta &= \beta_0 \times (1 + \text{energy}_{\text{ratio}})
\end{align}

où :
\begin{itemize}
    \item $\text{SNR}$ : Rapport signal/bruit (priorité à l'information si signal fort)
    \item $\text{confidence}_{\text{GP}}$ : Confiance du \gp{} (moins d'exploration si confiant)
    \item $\text{energy}_{\text{ratio}}$ : Ratio énergie consommée / énergie disponible
\end{itemize}

\subsection{Processus Gaussiens (\gp{}) pour l'Apprentissage Actif}

\subsubsection{Définition Mathématique}

Un \textbf{Processus Gaussien} est une distribution de probabilité sur les fonctions :

\begin{equation}
f(\mathbf{x}) \sim \mathcal{GP}(\mu(\mathbf{x}), k(\mathbf{x}, \mathbf{x}'))
\end{equation}

où :
\begin{itemize}
    \item $\mu(\mathbf{x}) = \mathbb{E}[f(\mathbf{x})]$ : Fonction moyenne (prédiction)
    \item $k(\mathbf{x}, \mathbf{x}') = \text{Cov}[f(\mathbf{x}), f(\mathbf{x}')]$ : Fonction de covariance (kernel)
\end{itemize}

\subsubsection{Kernel RBF (Radial Basis Function)}

Le kernel utilisé est une combinaison :

\begin{equation}
k(\mathbf{x}, \mathbf{x}') = \sigma^2 \times \exp\left(-\frac{\|\mathbf{x} - \mathbf{x}'\|^2}{2\ell^2}\right) + \sigma^2_{\text{noise}} \times \delta(\mathbf{x}, \mathbf{x}')
\end{equation}

où :
\begin{itemize}
    \item $\sigma^2$ : Variance du signal (amplitude)
    \item $\ell$ : Longueur d'échelle (corrélation spatiale)
    \item $\sigma^2_{\text{noise}}$ : Variance du bruit (incertitude de mesure)
    \item $\delta(\mathbf{x}, \mathbf{x}')$ : Fonction de Dirac (bruit blanc)
\end{itemize}

\textbf{Paramètres typiques :}
\begin{itemize}
    \item $\sigma^2 = 1.0$ : Variance du signal
    \item $\ell = 5.0$--$10.0$ m : Longueur d'échelle (corrélation sur 5--10 mètres)
    \item $\sigma^2_{\text{noise}} = 10^{-4}$ : Bruit de mesure très faible
\end{itemize}

\subsubsection{Prédiction \gp{}}

Étant donné $n$ observations $\mathcal{D} = \{(\mathbf{x}_i, y_i)\}_{i=1}^n$, la prédiction en un nouveau point $\mathbf{x}_*$ est :

\textbf{Moyenne (prédiction) :}
\begin{equation}
\mu(\mathbf{x}_*) = \mathbf{k}_*^T (\mathbf{K} + \sigma^2_{\text{noise}} \mathbf{I})^{-1} \mathbf{y}
\end{equation}

\textbf{Variance (incertitude) :}
\begin{equation}
\sigma^2(\mathbf{x}_*) = k(\mathbf{x}_*, \mathbf{x}_*) - \mathbf{k}_*^T (\mathbf{K} + \sigma^2_{\text{noise}} \mathbf{I})^{-1} \mathbf{k}_*
\end{equation}

où :
\begin{itemize}
    \item $\mathbf{K}_{ij} = k(\mathbf{x}_i, \mathbf{x}_j)$ : Matrice de covariance ($n \times n$)
    \item $\mathbf{k}_{*i} = k(\mathbf{x}_*, \mathbf{x}_i)$ : Vecteur de covariance ($n \times 1$)
    \item $\mathbf{y} = [y_1, \ldots, y_n]^T$ : Vecteur des observations
\end{itemize}

\textbf{Complexité computationnelle :}
\begin{itemize}
    \item Calcul de $\mathbf{K}$ : $\mathcal{O}(n^2)$
    \item Inversion de $(\mathbf{K} + \sigma^2_{\text{noise}} \mathbf{I})$ : $\mathcal{O}(n^3)$
    \item Prédiction : $\mathcal{O}(n^2)$ par point
\end{itemize}

\subsubsection{Fonction d'Acquisition UCB (Upper Confidence Bound)}

La fonction d'acquisition UCB équilibre exploration et exploitation :

\begin{equation}
\text{UCB}(\mathbf{x}) = \mu(\mathbf{x}) + \beta \cdot \sigma(\mathbf{x})
\end{equation}

où :
\begin{itemize}
    \item $\mu(\mathbf{x})$ : Prédiction (exploitation)
    \item $\sigma(\mathbf{x})$ : Incertitude (exploration)
    \item $\beta$ : Paramètre d'exploration (typiquement 2.0--3.0)
\end{itemize}

\textbf{Interprétation :}
\begin{itemize}
    \item Si $\beta$ est grand : Priorité à l'exploration (zones incertaines)
    \item Si $\beta$ est petit : Priorité à l'exploitation (zones à forte concentration)
\end{itemize}

\textbf{Sélection du point optimal :}
\begin{equation}
\mathbf{x}^* = \arg\max_{\mathbf{x} \in \mathcal{X}} \text{UCB}(\mathbf{x})
\end{equation}

où $\mathcal{X}$ est l'espace de recherche (grille $100 \times 100$ ou $150 \times 150$ selon le nombre d'observations).

La fonction d'acquisition utilise des poids dynamiques selon le nombre d'observations :
\begin{itemize}
    \item \textbf{Phase d'exploration} ($< 10$ observations) : 70\% incertitude, 30\% concentration
    \item \textbf{Phase équilibrée} ($10$--$50$ observations) : 50\% incertitude, 50\% concentration
    \item \textbf{Phase d'exploitation} ($> 50$ observations) : 30\% incertitude, 70\% concentration
\end{itemize}

\subsubsection{Stratégie Multi-Phase de Convergence}

Le \teacher{} utilise une stratégie adaptative selon la distance à la source estimée :

\textbf{Phase 1 : Exploration Active ($> 15$ m)}
\begin{itemize}
    \item Priorité à l'exploration (zones à haute incertitude)
    \item Pas de mouvement : 2--3.5 m
    \item Grille de recherche : $150 \times 150$ si $< 20$ observations, $100 \times 100$ sinon
\end{itemize}

\textbf{Phase 2 : Convergence Fine Guidée (5--15 m)}
\begin{itemize}
    \item Combinaison : 70\% vers source estimée + 30\% gradient
    \item Pas adaptatifs : 0.2--2.25 m (plus petit quand plus proche)
    \item Utilise l'estimation \gp{} pour guider la convergence
\end{itemize}

\textbf{Phase 3 : Recherche Locale en Spirale ($< 5$ m)}
\begin{itemize}
    \item Mouvement circulaire autour de la source estimée
    \item 60\% tangentiel (exploration) + 40\% radial (convergence)
    \item Pas très petits : 0.2--0.5 m
    \item Évite de dépasser la source
\end{itemize}

\subsection{Apprentissage par Renforcement (\rl{})}

\subsubsection{Problème de Contrôle Optimal}

Le \student{} apprend une politique $\pi_\theta(s) \to a$ qui maximise l'espérance de récompense cumulée :

\begin{equation}
J(\theta) = \mathbb{E}_{\tau \sim \pi_\theta} \left[\sum_{t=0}^T \gamma^t R(s_t, a_t)\right]
\end{equation}

où :
\begin{itemize}
    \item $\tau = (s_0, a_0, r_0, s_1, a_1, r_1, \ldots)$ : Trajectoire
    \item $\gamma \in [0, 1]$ : Facteur de discount (typiquement 0.99)
    \item $T$ : Horizon temporel
\end{itemize}

\subsubsection{Architecture du Réseau de Neurones}

Le \student{} utilise un réseau feedforward profond :

\begin{align}
\text{Input (État)} &: 16 \text{ dimensions} \\
\text{Couche 1} &: \text{Linear}(16 \to 256) + \text{Tanh} \\
\text{Couche 2} &: \text{Linear}(256 \to 256) + \text{Tanh} \\
\text{Couche 3} &: \text{Linear}(256 \to 128) + \text{Tanh} \\
\text{Output (Action)} &: \text{Linear}(128 \to 3) + \text{Tanh}
\end{align}

\textbf{Caractéristiques :}
\begin{itemize}
    \item Fonction d'activation : $\text{Tanh}$ pour borner les sorties dans $[-1, 1]$
    \item Initialisation : Xavier/Glorot uniform pour stabilité
\end{itemize}

\subsubsection{Algorithme DQN (Deep Q-Network)}

Le \student{} utilise une variante de DQN avec :
\begin{itemize}
    \item \textbf{Réseau principal} : $Q_\theta(s, a)$ -- Estimation de la valeur Q
    \item \textbf{Réseau cible} : $Q_{\theta'}(s, a)$ -- Cible stable pour l'apprentissage
    \item \textbf{Buffer d'expérience} : Stockage de $(s, a, r, s', \text{done})$
    \item \textbf{Mise à jour périodique} : $\theta' \leftarrow \theta$ toutes les 100 étapes
\end{itemize}

\textbf{Fonction de perte \rl{} :}
\begin{equation}
\mathcal{L}_{\text{RL}}(\theta) = \mathbb{E}\left[(Q_\theta(s, a) - y_{\text{target}})^2\right]
\end{equation}

où la cible est :
\begin{equation}
y_{\text{target}} = r + \gamma \cdot \max_{a'} Q_{\theta'}(s', a') \cdot (1 - \text{done})
\end{equation}

\subsubsection{Distillation de Connaissance}

Le \student{} apprend aussi du \teacher{} via la \textbf{divergence de Kullback-Leibler (KL)} :

\begin{equation}
\mathcal{L}_{\text{KL}}(\theta) = D_{\text{KL}}(\pi_{\text{teacher}} \| \pi_{\text{student}}) = \sum_a \pi_{\text{teacher}}(a|s) \cdot \log\frac{\pi_{\text{teacher}}(a|s)}{\pi_{\text{student}}(a|s)}
\end{equation}

\textbf{Politique du \teacher{} :}
\begin{equation}
\pi_{\text{teacher}}(a|s) = \text{softmax}\left(\frac{Q_{\text{teacher}}(s, a)}{T}\right)
\end{equation}

où $T$ est la température de distillation (typiquement 3.0).

\textbf{Politique du \student{} :}
\begin{equation}
\pi_{\text{student}}(a|s) = \text{softmax}\left(\frac{Q_\theta(s, a)}{T}\right)
\end{equation}

\textbf{Perte totale :}
\begin{equation}
\mathcal{L}(\theta) = \mathcal{L}_{\text{RL}}(\theta) + \lambda \cdot \mathcal{L}_{\text{KL}}(\theta)
\end{equation}

où $\lambda$ est le poids de distillation (typiquement 0.1--0.2).

\subsubsection{Exploration vs Exploitation}

Le \student{} utilise une stratégie \textbf{$\varepsilon$-greedy} :

\begin{equation}
a = \begin{cases}
\arg\max_a Q_\theta(s, a) & \text{avec probabilité } (1 - \varepsilon) \\
\text{action aléatoire} & \text{avec probabilité } \varepsilon
\end{cases}
\end{equation}

\textbf{Décroissance adaptative de $\varepsilon$ :}
\begin{equation}
\varepsilon = \max(\varepsilon_{\min}, \varepsilon_{\max} - (\varepsilon_{\max} - \varepsilon_{\min}) \times \frac{\text{step}}{\text{decay}_{\text{steps}}})
\end{equation}

Si la perte récente est faible ($< 0.1$), la décroissance est accélérée ($\times 1.5$).

\subsection{Validateur \gp{} pour Estimation de Position}

\subsubsection{Modélisation Probabiliste}

Le validateur \gp{} accumule les mesures et modélise la carte de concentration :

\begin{equation}
f(\mathbf{x}) \sim \mathcal{GP}(\mu(\mathbf{x}), k(\mathbf{x}, \mathbf{x}'))
\end{equation}

Avec les mêmes principes que le \teacher{} \gp{}.

\subsubsection{Estimation de Position de Fuite}

La position de fuite est estimée en combinant \textbf{concentration} et \textbf{confiance} :

\begin{equation}
\text{Score}(\mathbf{x}) = 0.7 \times \mu_{\text{normalized}}(\mathbf{x}) + 0.3 \times \text{confidence}(\mathbf{x})
\end{equation}

où :
\begin{itemize}
    \item $\mu_{\text{normalized}}(\mathbf{x})$ : Concentration normalisée $[0, 1]$
    \item $\text{confidence}(\mathbf{x}) = 1 - \sigma_{\text{relative}}(\mathbf{x})$ : Confiance (inverse de l'incertitude relative)
\end{itemize}

\textbf{Filtrage d'incertitude :}
\begin{equation}
\sigma_{\text{relative}}(\mathbf{x}) = \frac{\sigma(\mathbf{x})}{|\mu(\mathbf{x})| + \varepsilon}
\end{equation}

Les zones avec $\sigma_{\text{relative}} > 0.5$ sont pénalisées.

\textbf{Position estimée :}
\begin{equation}
\mathbf{x}_{\text{leak}} = \arg\max_{\mathbf{x} \in \mathcal{X}} \text{Score}(\mathbf{x}) \quad \text{tel que} \quad \text{Score}(\mathbf{x}) \geq \text{threshold}
\end{equation}

où $\text{threshold} = 0.95$ (probabilité minimale).

\subsubsection{Arrêt Automatique}

La simulation s'arrête automatiquement quand :

\begin{equation}
\text{confidence}(\mathbf{x}_{\text{leak}}) \geq 0.85
\end{equation}

Cela garantit une détection fiable tout en économisant l'énergie.

\subsection{Détecteur Amélioré Multi-Critères}

\subsubsection{Validation Multi-Critères}

Chaque détection est validée selon plusieurs critères :

\begin{enumerate}
    \item \textbf{Seuil de concentration} : $C_{\text{measured}} > \text{threshold}$
    \item \textbf{Distance à la source} : $\text{distance} < \text{min}_{\text{distance}}$
    \item \textbf{Gradient fort} : $\|\nabla C\| > \text{gradient}_{\text{threshold}}$
    \item \textbf{Confiance temporelle} : Détections cohérentes dans le temps
\end{enumerate}

\subsubsection{Estimation Robuste de Position}

L'estimation utilise \textbf{toutes les détections validées} avec :

\textbf{1. Clustering Spatial}
\begin{itemize}
    \item Identification du cluster principal (densité spatiale)
    \item Filtrage des outliers (distance $> 1.5 \times$ médiane du cluster)
\end{itemize}

\textbf{2. Pondération Temporelle}
\begin{itemize}
    \item Poids plus élevé pour les détections récentes
    \item Décroissance exponentielle : $w_t = \exp(-\lambda \cdot (t_{\max} - t))$
\end{itemize}

\textbf{3. Pondération de Cohérence Spatiale}
\begin{itemize}
    \item Poids plus élevé pour les détections proches d'autres détections
    \item Mesure de densité locale
\end{itemize}

\textbf{4. Estimation Finale}
\begin{itemize}
    \item Combinaison de moyenne pondérée et médiane pondérée
    \item Utilisation des 50 meilleures détections
\end{itemize}

\section{Architecture du Système}

\subsection{Vue d'Ensemble}

L'architecture du système \highlight{} est organisée en plusieurs couches interconnectées :

\begin{enumerate}
    \item \textbf{Interface Utilisateur} : Streamlit pour la configuration et la visualisation
    \item \textbf{Environnement de Simulation} : Gymnasium pour la simulation \rl{}
    \item \textbf{Capteur TDLAS} : Simulateur de capteur laser avec bruit réaliste
    \item \textbf{Modèle de Panache} : Simulation physique du panache de méthane
    \item \textbf{Architecture Teacher-Student} : \teacher{} (\gp{}) et \student{} (\rl{})
    \item \textbf{Détecteur Amélioré} : Validation multi-critères avec clustering
    \item \textbf{Validateur \gp{}} : Estimation probabiliste de position de fuite
    \item \textbf{Validateur de Performance} : Comparaison position réelle vs détectée
\end{enumerate}

\subsection{Composants Principaux}

\subsubsection{Modèle de Panache}

\textbf{Responsabilités :}
\begin{itemize}
    \item Calcul de la concentration $C(x, y, t)$
    \item Calcul du gradient $\nabla C(x, y, t)$
    \item Modélisation de l'advection par le vent
\end{itemize}

\subsubsection{Capteur TDLAS}

\textbf{Responsabilités :}
\begin{itemize}
    \item Simulation réaliste des mesures avec bruit
    \item Modélisation du rapport signal/bruit (SNR)
    \item Bruit atmosphérique et instrumental
\end{itemize}

\textbf{Modèle de mesure :}
\begin{equation}
C_{\text{measured}} = C_{\text{real}} \times (1 + \varepsilon_{\text{atmospheric}}) + \varepsilon_{\text{instrumental}}
\end{equation}

où :
\begin{itemize}
    \item $\varepsilon_{\text{atmospheric}} \sim \mathcal{N}(0, \sigma^2_{\text{atm}})$ : Bruit atmosphérique
    \item $\varepsilon_{\text{instrumental}} \sim \mathcal{N}(0, \sigma^2_{\text{inst}})$ : Bruit du capteur
\end{itemize}

\subsubsection{Environnement}

\textbf{Responsabilités :}
\begin{itemize}
    \item Implémentation de l'interface Gymnasium
    \item Gestion de la dynamique du drone
    \item Calcul des récompenses éco-informatives
    \item Intégration avec le modèle de panache et le capteur
\end{itemize}

\subsubsection{\teacher{} \gp{}}

\textbf{Responsabilités :}
\begin{itemize}
    \item Modélisation \gp{} de la carte de concentration
    \item Sélection des points d'exploration (acquisition UCB)
    \item Stratégie multi-phase de convergence
\end{itemize}

\subsubsection{\student{} \rl{}}

\textbf{Responsabilités :}
\begin{itemize}
    \item Apprentissage de la politique de navigation
    \item Distillation de connaissance depuis le \teacher{}
    \item Gestion du buffer d'expérience
\end{itemize}

\subsubsection{Validateur \gp{}}

\textbf{Responsabilités :}
\begin{itemize}
    \item Accumulation des mesures
    \item Estimation probabiliste de la position de fuite
    \item Calcul de la confiance
\end{itemize}

\subsubsection{Détecteur Amélioré}

\textbf{Responsabilités :}
\begin{itemize}
    \item Validation multi-critères des détections
    \item Estimation robuste de position (clustering, filtrage)
    \item Intégration avec le validateur \gp{}
\end{itemize}

\subsubsection{Validateur de Performance}

\textbf{Responsabilités :}
\begin{itemize}
    \item Comparaison position réelle vs détectée
    \item Calcul des métriques de performance
    \item Génération de rapports
\end{itemize}

\textbf{Métriques calculées :}
\begin{itemize}
    \item Taux de détection
    \item Précision de localisation (erreur, angle)
    \item Temps de détection
    \item Score global (0--100)
\end{itemize}

\subsection{Flux de Données}

Le flux de données suit les étapes suivantes :

\begin{enumerate}
    \item \textbf{Configuration Utilisateur} : Définition des paramètres (panache, capteur, drone, IA)
    \item \textbf{Initialisation Environnement} : Création des objets (PlumeConfig, TDLASConfig, EnvironmentConfig)
    \item \textbf{Initialisation IA} : Création du \teacher{}, \student{} (si mode full\_learning), et détecteur
    \item \textbf{Boucle de Simulation} (pour chaque étape) :
    \begin{itemize}
        \item Calcul concentration réelle via le modèle de panache
        \item Mesure capteur avec bruit réaliste
        \item Mise à jour \teacher{} avec nouvelle observation
        \item Sélection action (selon mode : \teacher{}, \student{}, ou combinaison)
        \item Exécution action dans l'environnement
        \item Validation détection par le détecteur amélioré
        \item Mise à jour validateur \gp{} avec nouvelle mesure
        \item Estimation position de fuite (priorité au validateur \gp{})
        \item Vérification arrêt automatique (si confiance $\geq 0.85$)
        \item Mise à jour \student{} (stockage expérience + apprentissage)
    \end{itemize}
    \item \textbf{Calcul Métriques Finales} : Comparaison position réelle vs détectée
    \item \textbf{Affichage Résultats} : Métriques, visualisations, rapport de validation
\end{enumerate}

\section{Implémentation Technique}

\subsection{Technologies Utilisées}

\begin{itemize}
    \item \textbf{Python 3.8+} : Langage principal
    \item \textbf{PyTorch} : Réseaux de neurones (\student{} \rl{})
    \item \textbf{scikit-learn} : Processus Gaussiens (\teacher{}, Validateur \gp{})
    \item \textbf{Gymnasium} : Environnement de simulation \rl{}
    \item \textbf{Streamlit} : Interface utilisateur
    \item \textbf{NumPy} : Calculs numériques
    \item \textbf{Plotly} : Visualisations interactives
    \item \textbf{Matplotlib} : Graphiques statiques
\end{itemize}

\subsection{Structure du Code}

La structure du projet est organisée en modules :

\begin{itemize}
    \item \texttt{models/} : \teacher{} \gp{} et \student{} \rl{}
    \item \texttt{simulation/} : Environnement Gymnasium et modèle de panache
    \item \texttt{sensors/} : Simulateur capteur TDLAS
    \item \texttt{analysis/} : Détecteur amélioré, validateur \gp{}, validateur de performance
    \item \texttt{experiments/} : Comparaisons expérimentales et tests de robustesse
\end{itemize}

\subsection{Détails d'Implémentation Clés}

\subsubsection{\teacher{} \gp{} -- Sélection de Point}

La méthode de sélection du prochain point à explorer implémente une stratégie multi-phase :

\begin{enumerate}
    \item \textbf{Calcul de la distance} à la source estimée (via \gp{} ou cible)
    \item \textbf{Phase 1} ($< 5$ m) : Recherche locale en spirale
    \begin{itemize}
        \item Mouvement circulaire autour de la source estimée
        \item Combinaison 60\% tangentiel + 40\% radial
        \item Pas de 0.2--0.5 m
    \end{itemize}
    \item \textbf{Phase 2} (5--15 m) : Convergence fine guidée
    \begin{itemize}
        \item Combinaison 70\% vers source + 30\% gradient
        \item Pas adaptatifs 0.2--2.25 m
    \end{itemize}
    \item \textbf{Phase 3} ($> 15$ m) : Exploration active
    \begin{itemize}
        \item Utilisation fonction d'acquisition UCB
        \item Grille adaptative selon nombre d'observations
        \item Sélection du maximum UCB
    \end{itemize}
    \item \textbf{Normalisation} et application du mouvement
\end{enumerate}

\subsubsection{\student{} \rl{} -- Apprentissage}

L'algorithme d'apprentissage combine perte \rl{} et distillation :

\begin{enumerate}
    \item \textbf{Échantillonnage} d'un batch depuis le buffer d'expérience
    \item \textbf{Calcul perte \rl{}} : Erreur quadratique entre $Q_\theta(s, a)$ et cible
    \item \textbf{Calcul perte KL} (si \teacher{} disponible) :
    \begin{itemize}
        \item Calcul politique \teacher{} et \student{} avec température
        \item Divergence KL entre les deux politiques
    \end{itemize}
    \item \textbf{Perte totale} : $\mathcal{L}_{\text{RL}} + \lambda \cdot \mathcal{L}_{\text{KL}}$
    \item \textbf{Rétropropagation} et mise à jour des paramètres
    \item \textbf{Clipping des gradients} pour stabilité
\end{enumerate}

\subsubsection{Validateur \gp{} -- Estimation de Position}

L'estimation de position suit une stratégie probabiliste :

\begin{enumerate}
    \item \textbf{Création grille adaptative} : $150 \times 150$ si $< 10$ mesures, $100 \times 100$ sinon
    \item \textbf{Prédiction \gp{}} : Calcul $\mu$ et $\sigma$ sur tous les points de la grille
    \item \textbf{Normalisation concentration} : $[0, 1]$
    \item \textbf{Calcul confiance} : Inverse de l'incertitude relative
    \item \textbf{Score combiné} : $0.7 \times \mu_{\text{normalized}} + 0.3 \times \text{confidence}$
    \item \textbf{Filtrage incertitude} : Pénalisation zones avec $\sigma_{\text{relative}} > 0.5$
    \item \textbf{Sélection position} : Maximum du score avec seuil adaptatif
\end{enumerate}

\section{Résultats et Validation}

\subsection{Métriques de Performance}

\subsubsection{Taux de Détection}

\begin{table}[H]
\centering
\begin{tabular}{lcc}
\toprule
\textbf{Mode} & \textbf{Taux de Détection} & \textbf{Amélioration vs Naïve} \\
\midrule
\teacher{} (\gp{}) & 85--92\% & +70\% à +77\% \\
\student{} (\rl{}) & 92--95\% & +77\% à +80\% \\
Full Learning & 93--96\% & +78\% à +81\% \\
Baseline Naïve & 12--15\% & -- \\
\bottomrule
\end{tabular}
\caption{Taux de détection par mode de simulation}
\end{table}

\textbf{Interprétation :}
\begin{itemize}
    \item Le \teacher{} explore intelligemment et détecte systématiquement
    \item Le \student{}, après apprentissage, dépasse même le \teacher{}
    \item Le mode Full Learning combine les deux approches pour un résultat optimal
\end{itemize}

\subsubsection{Précision de Localisation}

\begin{table}[H]
\centering
\begin{tabular}{lcc}
\toprule
\textbf{Mode} & \textbf{Erreur Moyenne} & \textbf{Distribution des Erreurs} \\
\midrule
\teacher{} (\gp{}) & 2.1 m & $< 2$ m : 60\%, 2--5 m : 30\%, $> 5$ m : 10\% \\
\student{} (\rl{}) & 1.8 m & $< 2$ m : 70\%, 2--5 m : 25\%, $> 5$ m : 5\% \\
Full Learning & 1.6 m & $< 2$ m : 75\%, 2--5 m : 20\%, $> 5$ m : 5\% \\
Baseline Naïve & $> 10$ m & $< 10$ m : 15\%, $> 10$ m : 85\% \\
\bottomrule
\end{tabular}
\caption{Précision de localisation par mode}
\end{table}

\textbf{Facteurs d'amélioration :}
\begin{enumerate}
    \item \textbf{Validateur \gp{}} : Estimation probabiliste précise
    \item \textbf{Stratégie multi-phase} : Convergence fine sans dépassement
    \item \textbf{Recherche locale en spirale} : Localisation précise $< 5$ m
\end{enumerate}

\subsubsection{Temps de Détection}

\begin{table}[H]
\centering
\begin{tabular}{lcc}
\toprule
\textbf{Mode} & \textbf{Temps Moyen} & \textbf{Amélioration} \\
\midrule
\teacher{} (\gp{}) & 2--12 s & -83\% \\
\student{} (\rl{}) & 0.8--2.5 s & -93\% \\
Full Learning & 0.7--2.0 s & -94\% \\
Baseline Naïve & 12.2 s & -- \\
\bottomrule
\end{tabular}
\caption{Temps de détection par mode}
\end{table}

\textbf{Explication :}
\begin{itemize}
    \item Le \student{}, après apprentissage, converge rapidement vers la source
    \item L'arrêt automatique (confiance $\geq 85\%$) économise du temps
\end{itemize}

\subsubsection{Efficacité Énergétique}

\begin{table}[H]
\centering
\begin{tabular}{lcc}
\toprule
\textbf{Mode} & \textbf{Efficacité (dét/kJ)} & \textbf{Économie d'Énergie} \\
\midrule
\teacher{} (\gp{}) & 0.15 & -25\% \\
\student{} (\rl{}) & 0.19 & -35\% \\
Full Learning & 0.20 & -37\% \\
Baseline Naïve & 0.08 & -- \\
\bottomrule
\end{tabular}
\caption{Efficacité énergétique par mode}
\end{table}

\textbf{Facteurs :}
\begin{itemize}
    \item Trajectoires optimisées (pas de zigzag inutile)
    \item Arrêt automatique (économie de 20--30\% d'énergie)
    \item Navigation guidée (moins de déplacements aléatoires)
\end{itemize}

\subsubsection{Score Global}

Le score global combine toutes les métriques :

\begin{equation}
\text{Score}_{\text{Global}} = 0.4 \times \text{Score}_{\text{Détection}} + 0.4 \times \text{Score}_{\text{Localisation}} + 0.2 \times \text{Score}_{\text{Efficacité}}
\end{equation}

\begin{table}[H]
\centering
\begin{tabular}{lcc}
\toprule
\textbf{Mode} & \textbf{Score Global} & \textbf{Interprétation} \\
\midrule
\teacher{} (\gp{}) & 70--85/100 & Bon à Excellent \\
\student{} (\rl{}) & 75--90/100 & Excellent \\
Full Learning & 80--92/100 & Excellent à Parfait \\
Baseline Naïve & 25--40/100 & Insuffisant \\
\bottomrule
\end{tabular}
\caption{Score global par mode}
\end{table}

\subsection{Validation Expérimentale}

\subsubsection{Protocole de Validation}

Le protocole de validation suit les étapes suivantes :

\begin{enumerate}
    \item \textbf{Configuration} : Position de fuite définie par l'utilisateur
    \item \textbf{Simulation} : Exécution avec le modèle \highlight{}
    \item \textbf{Détection} : Le système détecte automatiquement la position
    \item \textbf{Comparaison} : Calcul de l'erreur entre position réelle et détectée
    \item \textbf{Validation} : Vérification si erreur $< \text{tolérance}$ (10 m par défaut)
\end{enumerate}

\subsubsection{Résultats de Validation}

\textbf{Taux de succès mission} : 85--90\%

\begin{itemize}
    \item \textbf{Succès} : Erreur $\leq 10$ m (détection dans tolérance)
    \item \textbf{Partiel} : Erreur 10--20 m (détection proche)
    \item \textbf{Échec} : Erreur $> 20$ m (rare, $< 5\%$)
\end{itemize}

\textbf{Distribution des erreurs :}
\begin{itemize}
    \item $< 2$ m : 65\% des cas (excellente précision)
    \item 2--5 m : 25\% des cas (bonne précision)
    \item 5--10 m : 8\% des cas (précision acceptable)
    \item $> 10$ m : 2\% des cas (échecs)
\end{itemize}

\subsubsection{Robustesse}

Tests effectués sur \textbf{50+ positions différentes} :
\begin{itemize}
    \item Positions centrales (50, 50)
    \item Positions périphériques (10, 10), (90, 90)
    \item Positions aléatoires
    \item Grilles $3 \times 3$, $5 \times 5$
\end{itemize}

\textbf{Résultat} : Taux de succès $> 85\%$ sur toutes les positions testées.

\subsection{Comparaison avec l'État de l'Art}

\begin{table}[H]
\centering
\begin{tabular}{lccc}
\toprule
\textbf{Critère} & \textbf{\highlight{}} & \textbf{Méthodes Traditionnelles} & \textbf{Amélioration} \\
\midrule
Taux de détection & 92--95\% & 12--15\% & +600\% \\
Précision & 1.8 m & $> 10$ m & -82\% \\
Temps & 0.8--2.5 s & 12.2 s & -93\% \\
Énergie & Optimisée & Non optimisée & -35\% \\
Autonomie & Complète & Partielle & +100\% \\
\bottomrule
\end{tabular}
\caption{Comparaison avec l'état de l'art}
\end{table}

\section{Conclusion et Perspectives}

\subsection{Contributions Principales}

\begin{enumerate}
    \item \textbf{Architecture Teacher-Student innovante} : Combinaison \gp{} + \rl{} avec distillation
    \item \textbf{Fonction de récompense éco-informative} : Équilibre information/énergie
    \item \textbf{Validateur \gp{} probabiliste} : Estimation robuste avec arrêt automatique
    \item \textbf{Stratégie multi-phase} : Convergence précise sans dépassement
    \item \textbf{Détecteur multi-critères} : Validation robuste avec clustering
\end{enumerate}

\subsection{Résultats Clés}

\begin{itemize}
    \item \textbf{Taux de détection} : 92--95\% (vs 12--15\% baseline)
    \item \textbf{Précision} : 1.8 m d'erreur moyenne
    \item \textbf{Temps de détection} : 0.8--2.5 s (vs 12.2 s)
    \item \textbf{Économie d'énergie} : -35\%
    \item \textbf{Score global} : 80--92/100 (vs 25--40/100)
\end{itemize}

\subsection{Fiabilité Prouvée}

Le système inclut une \textbf{validation automatique} qui :
\begin{itemize}
    \item Compare position réelle (configurée) vs position détectée (indépendante)
    \item Calcule l'erreur de localisation
    \item Vérifie la tolérance (10 m)
    \item Génère des rapports détaillés
\end{itemize}

\textbf{Preuve de fiabilité} : Taux de succès 85--90\% sur positions variées.

\subsection{Perspectives}

\subsubsection{Phase 1 : Simulation (Actuel) ✓}
\begin{itemize}
    \item Validation de l'approche
    \item Optimisation des paramètres
    \item Tests de robustesse
\end{itemize}

\subsubsection{Phase 2 : Prototype Terrain (Prochaine étape)}
\begin{itemize}
    \item Intégration avec drone réel
    \item Tests en conditions réelles
    \item Calibration des capteurs
\end{itemize}

\subsubsection{Phase 3 : Déploiement}
\begin{itemize}
    \item Système opérationnel
    \item Intégration infrastructure
    \item Monitoring continu
\end{itemize}

\subsection{Impact Environnemental}

\begin{itemize}
    \item \textbf{Réduction des émissions} : Détection précoce des fuites
    \item \textbf{Efficacité énergétique} : Optimisation des trajectoires
    \item \textbf{Autonomie} : Système autonome sans intervention
    \item \textbf{Précision} : Localisation précise pour réparation ciblée
\end{itemize}

\section*{Annexes}

\subsection*{A. Références Mathématiques}

\begin{enumerate}
    \item \textbf{Processus Gaussiens} : Rasmussen \& Williams (2006) -- ``Gaussian Processes for Machine Learning''
    \item \textbf{Apprentissage Actif} : Krause et al. (2008) -- ``Near-Optimal Sensor Placements''
    \item \textbf{Apprentissage par Renforcement} : Sutton \& Barto (2018) -- ``Reinforcement Learning: An Introduction''
    \item \textbf{Distillation de Connaissance} : Hinton et al. (2015) -- ``Distilling the Knowledge in a Neural Network''
\end{enumerate}

\subsection*{B. Paramètres Optimaux (Concours)}

Les paramètres optimaux pour le concours sont disponibles dans les fichiers de configuration du projet.

\subsection*{C. Code Source}

Repository GitHub : \url{https://github.com/ElProfesormika/Projet_HIGHLIGHT-_Natran_-_UTT}

\vspace{2cm}

\begin{center}
\textbf{Équipe \highlight{}}\\
\textit{Concours Innovation Natran x Fondation UTT -- 2025}
\end{center}

\end{document}


